\documentclass[11pt]{article}

% ----------------------------------------------------------------------
%  Creo Advisors — AI Tools Implementation Plan (ENT 105)
%  Expands slide 11 ("AI tools") into work breakdown, tool stack,
%  effort/time estimates, and answers "where does the AI that writes
%  the customized email actually come from?"
%  Compile: pdflatex creo-ai-tools-implementation.tex  (run twice for ToC)
% ----------------------------------------------------------------------

\usepackage[utf8]{inputenc}
\usepackage[T1]{fontenc}
\usepackage{textcomp}
\usepackage[margin=1in]{geometry}
\usepackage{booktabs}
\usepackage{array}
\usepackage{longtable}
\usepackage{enumitem}
\usepackage{titlesec}
\usepackage{xcolor}
\usepackage{tcolorbox}
\usepackage[hidelinks]{hyperref}
\usepackage{xurl}
\newcolumntype{L}[1]{>{\raggedright\arraybackslash}p{#1}}

\definecolor{creonavy}{RGB}{17,38,73}
\definecolor{creoaccent}{RGB}{27,94,140}
\definecolor{creocyan}{RGB}{30,180,225}

\titleformat{\section}{\large\bfseries\color{creonavy}}{\thesection}{0.6em}{}
\titleformat{\subsection}{\normalsize\bfseries\color{creonavy!85!black}}{\thesubsection}{0.5em}{}
\titlespacing*{\section}{0pt}{1.6ex plus .3ex}{0.7ex}
\setlist[itemize]{leftmargin=1.3em,topsep=2pt,itemsep=2.5pt}
\setlist[enumerate]{leftmargin=1.6em,topsep=2pt,itemsep=3pt}
\setlength{\parskip}{4pt plus 1pt minus 1pt}
\setlength{\parindent}{0pt}
\setlength{\emergencystretch}{2em}

\newtcolorbox{bigidea}{colback=creonavy!6,colframe=creonavy,
  boxrule=0.8pt,arc=2pt,left=8pt,right=8pt,top=6pt,bottom=6pt}
\newtcolorbox{callout}{colback=creoaccent!8,colframe=creoaccent,
  boxrule=0.6pt,arc=2pt,left=8pt,right=8pt,top=5pt,bottom=5pt}
\newtcolorbox{caution}{colback=orange!7,colframe=orange!75!black,
  boxrule=0.5pt,arc=2pt,left=7pt,right=7pt,top=5pt,bottom=5pt}
\newtcolorbox{courseangle}{colback=green!5,colframe=green!45!black,
  boxrule=0.5pt,arc=2pt,left=7pt,right=7pt,top=5pt,bottom=5pt}

\newcommand{\hrs}[1]{\textbf{#1\,h}}
\newcommand{\verify}{\textcolor{orange!80!black}{\textsuperscript{\textbf{[confirm]}}}}

\title{\textbf{\color{creonavy} Creo Advisors --- AI Tools Implementation Plan}\\[3pt]
\large Turning slide 11 into a buildable system: work, tools, time, cost}
\author{\normalsize Prepared by the ENT 105 team for Creo Advisors}
\date{\normalsize 14 October 2026 --- input to the final setup proposal due 2 December 2026}

\begin{document}
\maketitle
\vspace{-1.2em}

\begin{bigidea}
\textbf{Bottom line.} Slide 11 names five tools but not the plumbing between them --- and in
particular it never says \emph{where the AI that writes each customized email actually runs}.
HubSpot does not contain Claude; it has its own built-in assistant (Breeze), which cannot be
trained on Rich's voice the way a Claude skill can. So the personalization row needs an
explicit bridge. We recommend building it in two stages: \textbf{(1)~a human-in-the-loop
setup} --- Rich (or an assistant) uses the Claude app with a saved \emph{``Rich's voice''}
skill and the HubSpot connector to draft emails in batches, review them, and push them back
into HubSpot; then, once the voice is proven, \textbf{(2)~a small automated pipeline} --- a
script that calls the Claude API for each contact and writes the draft into a HubSpot contact
field that a HubSpot sequence sends. Total build effort is roughly \textbf{125--195 hours over
7--10 weeks}, plus about \textbf{12--19 hours of Rich's own time}. Running on \textbf{HubSpot Starter},
with the Claude script doing the automation that HubSpot Professional would otherwise sell, keeps
year-one platform cost under about \$1{,}000 instead of about \$13{,}000 (\S\ref{sec:platform}).
The AI itself costs \textbf{tens of dollars a month}; SourceScrub is the remaining big license.
\end{bigidea}

\tableofcontents
\vspace{0.5em}\hrule\vspace{0.7em}

% ======================================================================
\section{The key question: where does the email-writing AI come from?}
% ======================================================================
The slide says ``Claude for personalization \ldots\ plugs into HubSpot and email.'' Nothing
plugs in by default. There are three realistic ways to make that happen. They differ in who
presses the button, what it costs, and how much engineering is needed.

\begin{center}
\renewcommand{\arraystretch}{1.3}\small
\begin{tabular}{L{2.6cm} L{4.3cm} L{3.6cm} L{3.6cm}}
\toprule
\textbf{Option} & \textbf{How it works} & \textbf{Pros} & \textbf{Cons} \\
\midrule
\textbf{A. HubSpot Breeze} (HubSpot's built-in AI) &
HubSpot's own AI writing assistant, used inside the email editor and in sequences. &
Zero integration work; already in the HubSpot UI. &
Not Claude; it can't load Rich's voice skill or content bank. Output tends to sound generic
--- the exact drift problem Creo already saw. \\
\addlinespace
\textbf{B. Claude app + skill + HubSpot connector} \emph{(recommended start)} &
A Claude Team/Pro account holds a ``Rich's voice'' \emph{skill} and a Project with his content
bank. The user connects HubSpot, asks Claude to pull a list (e.g.\ ``ACG Boston attendees,
persona = Operating Partner''), drafts one email per contact, reviews them, and saves them back
as HubSpot drafts, tasks, or notes. &
Uses Claude and the voice skill directly; no code; a human reviews every send; live in about
2 weeks. &
Someone has to run it (about 30--60 min per batch); throughput tops out around 50--150 contacts
per session. \\
\addlinespace
\textbf{C. Claude API pipeline} \emph{(scale stage)} &
A small script, scheduled or run on demand, reads contacts from the HubSpot API, sends each one
plus the voice skill to the Claude API, and writes the draft into a custom contact property
(\texttt{ai\_email\_body}). A HubSpot sequence or workflow inserts that property as a
personalization token and sends it. &
Fully automated at any volume; cheap per email; keeps HubSpot as the system of record. &
Needs a developer for about 35--55 h and an API key; a review step must be designed in; the
script has to be maintained. \\
\bottomrule
\end{tabular}
\end{center}

\begin{callout}
\textbf{Recommendation.} Start with \textbf{Option B} so Rich can approve the voice on real
emails. Graduate to \textbf{Option C} once roughly 90\% of drafts go out with light or no
edits. Use \textbf{Option A} (Breeze) only for small things like subject-line suggestions.
Throughout, the \emph{same} voice skill is the single source of truth: the Claude app loads it
directly, and the API pipeline sends the same text as its system prompt.
\end{callout}

\subsection{What ``Claude Code + skills'' should mean on the slide}
Claude Code is a developer tool that runs in a terminal. Rich will not use it day to day. The
\emph{skills} it refers to are also available in the regular Claude desktop and web apps on
paid plans, and that is where Rich and his assistant will use them. Claude Code is
what the \emph{implementer} uses to write the skill files, build the Option~C script, and
automate reporting. We suggest the slide row read \emph{``Claude + skills (built with Claude
Code).''}

\subsection{The data flow, end to end}
\begin{center}
\small
\fbox{\parbox{0.92\linewidth}{\centering
\textbf{SourceScrub / ACG list} $\rightarrow$ CSV import $\rightarrow$
\textbf{HubSpot} (contact + persona + sector + company + event tag)\\[3pt]
$\downarrow$\\[1pt]
\textbf{Claude} (voice skill + content bank + contact record) $\rightarrow$ draft email /
LinkedIn note\\[3pt]
$\downarrow$\\[1pt]
\textbf{Human review} (Rich or assistant) $\rightarrow$ \textbf{HubSpot sequence / marketing
email} sends from Rich's mailbox\\[3pt]
$\downarrow$\\[1pt]
Opens, clicks, replies, meetings logged in \textbf{HubSpot} $\rightarrow$ \textbf{Claude for
analysis} (weekly summary by persona) $\rightarrow$ next post or next email}}
\end{center}
LinkedIn \emph{messages} cannot be automated through any legitimate API, so Claude drafts them
and a person sends them by hand. LinkedIn \emph{posts} can be scheduled through HubSpot's social
tool or LinkedIn's own scheduler.

% ======================================================================
\section{Tool stack}
% ======================================================================
\begin{center}
\renewcommand{\arraystretch}{1.25}\small
\begin{tabular}{L{3.4cm} L{5.6cm} L{4.6cm}}
\toprule
\textbf{Tool} & \textbf{Role} & \textbf{Tier / notes} \\
\midrule
\textbf{Claude app} (Team or Pro) & Voice skill, content-bank Project, HubSpot connector,
batch drafting, weekly analysis & Team plan if 2+ users need to share the skill and Project \\
\textbf{Claude API} & Option C pipeline; weekly reporting job &
Pay per use; Claude Opus 5.5 by default (Sonnet 5.5 if cost ever matters); Batch API halves the
cost for overnight runs \\
\textbf{Claude Code} & Implementer's tool to write skills, the sync script, and the reports &
Developer use only \\
\textbf{HubSpot Starter Customer Platform} & CRM, contacts, newsletter, marketing email,
sequences sent from Rich's own inbox, tracking, dashboards & \textbf{Recommended tier}, about
\$15--20 per seat per month; 2 seats (Rich + assistant). Upgrade to Professional only if volume
outgrows the script (\S\ref{sec:platform})\verify \\
\textbf{LinkedIn native scheduler} & Schedule Rich's posts & Free; replaces HubSpot
Professional's social publishing \\
\textbf{SourceScrub} & Conference and firm database; pre-ACG target lists & Annual enterprise
contract, quote-based\verify. Confirm it covers \emph{attendee} lists for ACG events, not only
sponsors and exhibitors \\
\textbf{Email domain (DNS)} & SPF, DKIM, DMARC so HubSpot sends land in the inbox & One-time
setup; required before any volume \\
\textbf{Small host for the script} & Runs the Option C script (GitHub Actions, a cloud
function, or a laptop cron job) & About \$0--10/month \\
\bottomrule
\end{tabular}
\end{center}

% ======================================================================
\section{Work breakdown by workstream}
% ======================================================================
Hours are for the implementer (a student team member, a contractor, or Creo staff). Rich's
time is listed separately because it is the scarcest input.

\subsection{WS1 --- Discovery \& decisions (week 1)}
\begin{itemize}
  \item Inventory where contacts live today (Outlook, Excel, LinkedIn export, phone) and how many
  there are; list past clients and their relationship owners.
  \item Decide HubSpot tiers, Claude plan, and whether SourceScrub is worth its contract versus
  ACG's own attendee app and LinkedIn Sales Navigator.
  \item Define the 3--4 personas (e.g.\ PE Operating Partner, Deal Partner, Portfolio CEO/CFO,
  Board member) and the sectors used for tagging.
  \item Write the review rule: who approves AI drafts before anything is sent.
\end{itemize}
\textbf{Effort:} \hrs{8--12} implementer; \hrs{2--3} Rich.

\subsection{WS2 --- Voice content bank \& ``Rich's voice'' skill (weeks 1--3)}
\begin{itemize}
  \item Collect 30--60 samples of Rich's writing: LinkedIn posts, sent emails he's proud of,
  deck narratives, case write-ups, and a short recorded interview transcript if written
  samples are thin.
  \item Distill them into a skill: tone rules, signature phrases, words he never uses, how he
  opens and closes, length limits, compliance rules (no client names without permission, no
  return claims), and 8--12 annotated example emails and posts by persona.
  \item Build a Claude Project holding the content bank (case summaries, service lines,
  proof points) so drafts can cite real material.
  \item \textbf{Calibration loop:} generate 20 drafts across personas; Rich marks each one ``sounds
  like me / doesn't''; revise the skill; repeat 2--3 rounds.
  \item Version the skill files (in Git or a shared folder) so the app and the API use the
  identical text. This is the fix for the long-prompt drift the slide mentions.
\end{itemize}
\textbf{Effort:} \hrs{20--30} implementer; \hrs{5--8} Rich (samples plus 2--3 review rounds).

\subsection{WS3 --- HubSpot foundation (weeks 1--4)}
\begin{itemize}
  \item Set up the account, users, and email-domain authentication (SPF/DKIM/DMARC); connect
  Rich's inbox and calendar.
  \item Create custom properties: \texttt{persona}, \texttt{sector}, \texttt{event\_source}
  (e.g.\ ACG Boston 2027), \texttt{relationship\_owner}, \texttt{past\_client},
  \texttt{ai\_email\_subject}, \texttt{ai\_email\_body}, \texttt{ai\_draft\_status}.
  \item Clean, de-duplicate, and import contacts and past clients; set lifecycle stages; record
  opt-in status so the newsletter is compliant (CAN-SPAM; consent for any non-U.S.\ contacts).
  \item Build the newsletter template and a monthly send; build 2--3 sequences (pre-event,
  post-event, past-client re-engagement).
  \item Set up post scheduling in LinkedIn's own scheduler (free); HubSpot social publishing
  needs the Professional tier.
  \item Build a dashboard: sends, opens, clicks, replies, and meetings, broken out by persona and
  by event.
\end{itemize}
\textbf{Effort:} \hrs{25--40} implementer (import cleanup is the wildcard); \hrs{2--3} Rich.

\subsection{WS4 --- Personalization engine (weeks 3--7)}
\textbf{Stage B (Claude app, weeks 3--4).}
\begin{itemize}
  \item Connect the HubSpot connector in Claude; test read access to lists and write access to
  notes, tasks, and properties\verify.
  \item Write a second skill, ``Outreach drafter,'' that takes a contact record plus a trigger
  (event, news, past engagement) and returns a subject, a 90--150-word body, and a LinkedIn
  note, all in Rich's voice. It must use only facts present in the record or the content bank
  and must never invent a personal detail.
  \item Document the operator runbook: pull list $\rightarrow$ draft $\rightarrow$ review
  $\rightarrow$ save to HubSpot $\rightarrow$ enroll in sequence.
\end{itemize}
\textbf{Stage C (API pipeline, weeks 5--7).}
\begin{itemize}
  \item Write a script (Python, built with Claude Code) that pulls contacts where
  \texttt{ai\_draft\_status = needed}, calls the Claude API with the voice skill as a cached
  system prompt, and writes subject, body, and \texttt{status = review} back to HubSpot.
  \item Set up a HubSpot view or task queue so the reviewer can approve
  (\texttt{status = approved}); the script then enrolls approved contacts in the sequence
  through the HubSpot API, or the reviewer enrolls them in a few clicks from a saved view if
  Starter's API does not allow enrollment\verify. The sequence inserts the
  \texttt{ai\_email\_body} token.
  \item Have the script take over the jobs HubSpot Professional workflows would otherwise do:
  updating lifecycle stages, tagging by event, creating follow-up tasks when someone replies or
  goes quiet, and flagging stale contacts. This is what lets Creo stay on Starter.
  \item Add guardrails: skip contacts missing a persona or company; log every draft; cap the
  batch size; rotate keys and keep them in a secrets store.
  \item Optional: enrich each contact with recent company news before drafting (Claude's
  web-search tool), with a cited source the reviewer can check.
\end{itemize}
\textbf{Effort:} Stage B \hrs{12--18}; Stage C \hrs{35--55}; \hrs{2--3} Rich (approving the first
batches).

\subsection{WS5 --- SourceScrub \& event prospecting (weeks 4--6, before the first ACG event)}
\begin{itemize}
  \item Run a demo or trial to confirm SourceScrub has usable attendee data for the specific ACG
  events Creo attends\verify.
  \item Build saved searches: PE firms by fund size and sector, and portfolio companies matching
  Creo's ideal client.
  \item Export, map fields to HubSpot properties, tag \texttt{event\_source}, and import.
  \item Run the pre-event sequence 2--3 weeks out, ask for a meeting at the event, and run a
  post-event follow-up within 48 hours.
\end{itemize}
\textbf{Effort:} \hrs{10--15} first event, then about \hrs{3--5} per event.

\subsection{WS6 --- Claude for analysis (weeks 6--8)}
\begin{itemize}
  \item A weekly report, generated by Claude through the HubSpot connector or by a scheduled API
  job reading a HubSpot export: results by persona and sector, which subject lines and topics
  worked, overdue follow-ups, and replies that need Rich.
  \item A monthly ``themes'' memo: recurring questions and objections from replies, which feed
  the next LinkedIn post (closing the loop the slide describes).
\end{itemize}
\textbf{Effort:} \hrs{8--14} implementer.

\subsection{WS7 --- Training, documentation \& handoff (weeks 7--9)}
\begin{itemize}
  \item Runbooks for each workflow; a 60-minute training session for Rich and his assistant; a
  one-page ``how to update the voice skill.''
  \item A 30-day hypercare period: fix issues and tune the skill based on reply data.
\end{itemize}
\textbf{Effort:} \hrs{8--12} implementer; \hrs{1--2} Rich.

% ======================================================================
\section{Time summary \& timeline}
% ======================================================================
\begin{center}
\renewcommand{\arraystretch}{1.2}\small
\begin{tabular}{L{6.4cm} r r L{2.6cm}}
\toprule
\textbf{Workstream} & \textbf{Implementer} & \textbf{Rich} & \textbf{Weeks} \\
\midrule
WS1 Discovery \& decisions & 8--12 h & 2--3 h & 1 \\
WS2 Voice bank \& skill & 20--30 h & 5--8 h & 1--3 \\
WS3 HubSpot foundation & 25--40 h & 2--3 h & 1--4 \\
WS4 Personalization (B then C) & 47--73 h & 2--3 h & 3--7 \\
WS5 SourceScrub \& event prospecting & 10--15 h & --- & 4--6 \\
WS6 Analysis \& reporting & 8--14 h & --- & 6--8 \\
WS7 Training \& handoff & 8--12 h & 1--2 h & 7--9 \\
\midrule
\textbf{Total} & \textbf{126--196 h} & \textbf{12--19 h} & \textbf{7--10 weeks} \\
\bottomrule
\end{tabular}
\end{center}
With one implementer at about 20 h/week this is roughly 7--10 calendar weeks. If Option C is
deferred, the total drops to about \textbf{90--140 h} and \textbf{5--6 weeks}. The critical
path runs through the voice calibration (WS2): nothing should go to contacts until Rich has
approved the voice.

\textbf{Milestones.}
\begin{enumerate}
  \item Week 2: HubSpot live and authenticated; first voice skill draft.
  \item Week 4: Voice approved; first human-reviewed batch sent with Option B.
  \item Week 6: First event campaign run end to end (SourceScrub $\rightarrow$ HubSpot
  $\rightarrow$ Claude $\rightarrow$ sequence).
  \item Week 8: Option C pipeline live with an approval step; weekly analysis report running.
  \item Week 9: Handoff and runbooks; begin 30-day hypercare.
\end{enumerate}

% ======================================================================
\section{Platform choice: HubSpot Professional vs.\ Starter vs.\ Klaviyo}\label{sec:platform}
% ======================================================================
We compared the email/CRM platform options at Creo's likely size (about 1{,}000--1{,}500
contacts). Prices are 2026 list prices reported by third-party pricing guides, not quotes for
Creo; the guides disagree on some bands, so confirm on each vendor's live pricing page.

\begin{center}
\renewcommand{\arraystretch}{1.3}\small
\begin{tabular}{L{3.6cm} L{2.6cm} L{2.4cm} L{5.2cm}}
\toprule
\textbf{Option} & \textbf{Approx.\ monthly} & \textbf{Year 1} & \textbf{Fit for Creo} \\
\midrule
HubSpot Marketing Hub Professional & \$800--890, 3 seats, 2{,}000 contacts &
\$13--14k incl.\ required \$3{,}000 onboarding & Full fit, but most of the price buys workflow
automation and social publishing that the Claude script and LinkedIn can cover \\
\addlinespace
Klaviyo Email & \$30--45 (501--1{,}500 active profiles) & \$400--550 \emph{plus a separate CRM} &
Poor fit: built for e-commerce, so it lacks a CRM and inbox sequences \\
\addlinespace
\textbf{HubSpot Starter + Claude script} \emph{(recommended)} & \$30--40 (2 seats at
\$15--20) & \textbf{\$400--500} & Good fit: one system of record and personal sequences, with
the automation handled by the script \\
\bottomrule
\end{tabular}
\end{center}

\subsection{Why not Klaviyo}
Klaviyo is cheaper on paper, but it is built for online retailers selling to consumers:
newsletters and SMS triggered by browsing and purchase events, billed per active profile. Creo
needs relationship outreach between businesses, and Klaviyo is missing the core pieces:
\begin{itemize}
  \item \textbf{No CRM.} It has no deals, pipeline, relationship owners, or meeting logging. Creo
  would still need to buy and connect a CRM, which undoes most of the saving and breaks the
  slide's ``one home for contacts, sends and results.''
  \item \textbf{No sales sequences from Rich's own inbox.} Klaviyo sends through a marketing
  setup, so a personal note to a PE partner would look like a mass email, which defeats the
  ``feel 1:1'' goal.
  \item \textbf{No LinkedIn scheduling and weak sales reporting.} The weekly Claude analysis would
  have to stitch together two systems.
\end{itemize}

\subsection{The real cost lever: Starter plus the script}
Most of HubSpot Professional's price pays for automated workflows and social publishing. In
this plan the Option~C script already does the hard part: selecting contacts, drafting,
tracking approval status, and enrolling contacts. Splitting the work this way keeps Creo on Starter:
\begin{itemize}
  \item \textbf{HubSpot Starter:} contacts, CRM, newsletter, basic sequences, open/click/reply
  tracking, dashboards.
  \item \textbf{Claude script:} drafting, approval status, sequence enrollment, lifecycle updates,
  follow-up tasks.
  \item \textbf{LinkedIn's own scheduler:} posts (free).
\end{itemize}
The trade-off is about \textbf{10--15 more build hours} (already included in WS4 and the time
summary) and one more script to maintain. Year-one platform cost falls from roughly \$13{,}000
to under \$1{,}000.

\begin{caution}
\textbf{Confirm with HubSpot before committing:} (1)~that Starter includes sequences at the
volume Creo plans to send; (2)~that Starter's API lets the script write the custom
\texttt{ai\_*} fields and enroll contacts in sequences; (3)~current Starter seat pricing,
including any new-customer promotion. \textbf{Upgrade path:} move to Professional later only if
outreach volume or reporting needs outgrow the script; the data model and skills carry over
unchanged.
\end{caution}

\textbf{Pricing sources (accessed 14 October 2026).}
{\small
\begin{itemize}
  \item SendX, Klaviyo pricing 2026: \url{https://www.sendx.io/blog/klaviyo-pricing-plans-costs-alternatives-2026}
  \item Spoks, Klaviyo pricing breakdown: \url{https://spoks.com/resources/klaviyo-pricing-breakdown}
  \item Omnisend, Klaviyo pricing: \url{https://www.omnisend.com/blog/klaviyo-pricing/}
  \item MoEngage, Klaviyo pricing guide 2026: \url{https://www.moengage.com/blog/klaviyo-pricing/}
  \item Tinycommand, HubSpot pricing explained (2026): \url{https://tinycommand.com/blogs/hubspot-pricing-explained}
  \item Agiled, HubSpot pricing 2026: \url{https://agiled.app/blog/hubspot-pricing}
  \item Featurebase, HubSpot pricing 2026: \url{https://featurebase.app/blog/hubspot-pricing}
  \item Automation Atlas, HubSpot pricing 2026: \url{https://automationatlas.io/answers/hubspot-pricing-explained-2026/}
\end{itemize}}

% ======================================================================
\section{Running cost (estimated)}
% ======================================================================
\begin{center}
\renewcommand{\arraystretch}{1.2}\small
\begin{tabular}{L{6.6cm} L{6.4cm}}
\toprule
\textbf{Item} & \textbf{Estimate} \\
\midrule
Claude app (Pro, or Team for shared skills) & About \$20--30 per user per month\verify \\
Claude API (Opus 5.5: \$4 / \$20 per million input/output tokens) &
About \$0.02--0.05 per drafted email; 1{,}000 emails/month is about \$20--50, and about half
that via Batch API with a cached voice prompt \\
HubSpot Starter Customer Platform, 2 seats & About \$30--40 per month
(about \$360--480 per year)\verify \\
SourceScrub & Annual contract; quote-based, often five figures\verify \\
Script hosting & \$0--10 per month \\
\bottomrule
\end{tabular}
\end{center}
With HubSpot on Starter, the platform and AI together run under about \$100 a month. The real
costs are the SourceScrub contract and the build hours. Prices marked \verify{} are list-price ranges we have not
quoted for Creo; confirm them with each vendor before the 12/2 proposal.

% ======================================================================
\section{Risks \& safeguards}
% ======================================================================
\begin{caution}
\begin{itemize}
  \item \textbf{Hallucinated personalization.} An invented ``congrats on the new role'' is worse
  than no personalization. The skill must use only facts in the record or a cited source, and a
  human reviews drafts until trust is earned.
  \item \textbf{Voice drift.} Keep one versioned skill file; re-run the 20-draft calibration
  whenever it changes.
  \item \textbf{Deliverability.} Authenticate the domain, warm up volume gradually, send 1:1
  outreach as sequences from Rich's inbox, and keep the newsletter on marketing email with
  unsubscribe handling.
  \item \textbf{Confidential client data.} Never put client deal data into prompts beyond what's
  already approved for marketing. Use a business Claude plan (commercial terms; no training on
  Creo's data) and keep API keys out of shared documents.
  \item \textbf{LinkedIn terms.} No automation tools that send DMs or scrape LinkedIn; drafts only.
  \item \textbf{SourceScrub fit.} Attendee coverage for specific ACG chapters is unverified ---
  run a trial before signing an annual contract.
\end{itemize}
\end{caution}

% ======================================================================
\section{Open questions for Creo}
% ======================================================================
\begin{enumerate}
  \item Where do contacts live today, and roughly how many are there?
  \item Who besides Rich will run the workflows (an assistant, an associate, or us)?
  \item Does Creo already have any HubSpot or SourceScrub licenses?
  \item Which ACG events in the next six months are the targets?
  \item What monthly outreach volume is the goal: 50, 500, or 5{,}000 contacts?
  \item Which past posts, decks, and cases can go into the content bank, and which are
  confidential?
\end{enumerate}

\begin{courseangle}
\textbf{Course angle (ENT 105).} This is the inbound engine the syllabus is built around:
HubSpot as the system of record, persona-driven messaging, and lead generation measured weekly.
The weekly Claude analysis report is designed to be the agenda for a weekly marketing management
meeting --- what we sent, to which persona, what worked, and what to do next week.
\end{courseangle}

\end{document}
